\documentclass[10pt,a4paper]{amsart}
\usepackage[margin=19mm]{geometry}
\usepackage{amsmath,amssymb,mathtools}
\usepackage[scr=rsfs,cal=euler]{mathalfa}
\usepackage{xfrac}
\usepackage[unicode,hidelinks,bookmarks=false]{hyperref}
\newcommand{\ie}{i.e.\ }
\newcommand{\cf}{cf.\ }
\newcommand{\resp}{respectively\ }
\newcommand{\Subsection}{\subsection}
\providecommand{\nobreakdash}{\nobreak\hbox{-}}
\setlength{\emergencystretch}{2em}
\multlinegap0pt
% ---------------------------------------------------------------------------
% Editorial AMS wrapper prelude. The theorem environments and the mathematical macros below are
% transcribed from the paper's own preamble (cached-originals/source0.tex lines 42--70); the AMS amsart
% class, the named format packages of the pinned TeX Live runtime and this editorial formatting are not
% part of the author's text. No mathematics is changed.
% ---------------------------------------------------------------------------
\newtheorem{axiom}{Axiom}
\newtheorem{claim}[axiom]{Claim}
\newtheorem{theorem}{Theorem}[section]
\newtheorem{lemma}[theorem]{Lemma}
\newtheorem{corollary}[theorem]{Corollary}
\newtheorem{proposition}[theorem]{Proposition}

\theoremstyle{definition}
\newtheorem{definition}[theorem]{Definition}
\newtheorem{example}[theorem]{Example}
\newtheorem{remark}[theorem]{Remark}
\newtheorem{assumption}[theorem]{Assumption}
\newtheorem*{fact}{Fact}

\def\R{{\mathbb R}}
\def\E{{\mathbb E}}
\def\P{{\mathbb P}}
\def\Z{{\mathbb Z}}


\begin{document}
\title{Cluster size distributions of discrete random fields}
\author{Dan Cheng and John Ginos}
\date{Source: arXiv:2601.14586v2; distributed under CC BY 4.0}
\maketitle
\noindent\emph{Source, version and licence.} \emph{Cluster size distributions of discrete random fields}, by Dan Cheng and John Ginos. The source of this editable unit is the e-print of \textbf{version v2} of arXiv:2601.14586, https://arxiv.org/abs/2601.14586v2, whose material is distributed under the Creative Commons Attribution 4.0 International licence, http://creativecommons.org/licenses/by/4.0/, as recorded in the arXiv licence metadata for this paper and shown on that abstract page. The whole of the body below is version v2, the newest version of the paper. Full rights notice: licenses/arxiv-2601.14586-CC-BY-4.0.txt.\par\smallskip
\noindent\textit{Editorial scope.} Counted unit: original Theorem 2.2 of the article, the statement carried by the source environment \texttt{theorem} with source label \texttt{thm:CSD\_d}, cached-originals/source0.tex lines 418--429, printed as \textbf{Theorem 2.2} exactly as in the article. It is delivered together with the authors' own complete proof as the single primary proof body, source0.tex lines 431--457 (27 lines): the \texttt{proof} environment that immediately follows the statement and derives the distribution formula from the defining identities. No proof marker is added and no mathematics is written, repaired or re-derived; the authors' notation and the printed slips of the source are kept literal. The closure of the counted statement and proof over references and citations was computed: it contains no citation at all, and the referenced definitions, the finite-cluster assumption and the weight formula that the proof uses are included in the delivered file as uncounted context, so every reference inside the retained passages resolves inside it. The article's own numbering is reproduced by explicit counter settings: the counter \texttt{theorem}, declared by \texttt{\textbackslash newtheorem\{theorem\}\{Theorem\}[section]} with lemma, corollary, proposition, definition, example, remark and assumption sharing it, is set before each retained passage, so the counted statement prints as Theorem 2.2 and the retained context keeps its original numbers, exactly as in the article. No \texttt{\textbackslash cite} command occurs in any retained passage, so this unit delivers no bibliography entry even though the article itself has a \texttt{thebibliography}; the bundled source retains it. The retained passages contain no figure environment and no graphics-inclusion command, so no artwork is redistributed and none of the many artwork files shipped by the e-print is used. The source's own class is the publisher class \texttt{imsart} of the Institute of Mathematical Statistics (\texttt{\textbackslash documentclass[aap,preprint]\{imsart\}}); that class and its style file are shipped by the e-print but are \emph{not} shipped, loaded or used by this package and are never redistributed, and the wrapper is the AMS \texttt{amsart} class with the article's theorem declarations restated and its macros transcribed from its preamble. Every declared body of this unit is an exact, unmodified span of source0.tex: no body is a bounded passage comparison, \texttt{omit} was not used anywhere, and consequently no fidelity receipt is asserted in delivery-evidence.json. The complete concatenated original source is bundled as sources/193098/source0.tex. Omitted from this editable unit and therefore absent from it are source0.tex lines 1--148; 157--394; 403--410; 417--417; 430--430; 458--2801, that is the article's own preamble, title and author block and abstract, the introduction, the simulation and empirical sections with their figures, the peak-based cluster size results and the Gaussian large-deviation theorems with their proofs, and the article's own bibliography entries, all of which remain present in the bundled source but are not delivered here. The AMS \texttt{amsart} wrapper, the named format packages of the pinned TeX Live runtime and this editorial source, licence and scope paragraph are editorial formatting only.\par\smallskip
\setcounter{section}{1}
\setcounter{equation}{0}
\setcounter{axiom}{0}
\setcounter{theorem}{0}
	\begin{assumption}[Finite clusters]\label{ass:finite-cluster}
		At the threshold $u$ under consideration,
		\begin{equation*}
			\mathbb P\left(
			|C_u(t)|<\infty \text{ for every }t\in E_u
			\right)=1.
		\end{equation*}
	\end{assumption}

\setcounter{section}{2}
\setcounter{equation}{6}
\setcounter{axiom}{0}
\setcounter{theorem}{1}
	\begin{equation}\label{eq:w_d}
		\begin{split}
			w_k
			&:=\sum_{D\in\mathcal{C}_k^{\rm root}}
			\P\bigl(X_D>u,\ X_{\mathcal{N}(D)}\le u\bigr),
			\qquad k=1,2,\ldots.
		\end{split}
	\end{equation}

\setcounter{section}{2}
\setcounter{equation}{7}
\setcounter{axiom}{0}
\setcounter{theorem}{1}
	\begin{equation}\label{eq:w_d_2}
		\begin{split}
			w_k =\P\bigl(R_u(o),\ |C_u(o)|=k\bigr),
			\qquad k=1,2,\ldots.
		\end{split}
	\end{equation}

\setcounter{section}{2}
\setcounter{equation}{8}
\setcounter{axiom}{0}
\setcounter{theorem}{1}
	\begin{theorem}\label{thm:CSD_d}
		Let $\{X_t:t\in\Z^d\}$ be stationary, let $u\in\R$ satisfy
		Assumption~\ref{ass:finite-cluster}, and suppose
		$\rho_u=\P(R_u(o))>0$.  Then
		\begin{equation}\label{eq:CSD_d}
			\P(S_u=k)=\frac{w_k}{\rho_u},
			\qquad
			\rho_u=\sum_{j\geq1}w_j,
			\qquad k\geq1,
		\end{equation}
		where $w_k$ is given by \eqref{eq:w_d}.
	\end{theorem}

\setcounter{section}{2}
\setcounter{equation}{9}
\setcounter{axiom}{0}
\setcounter{theorem}{2}
	\begin{proof}
		By stationarity, without loss of generality, we assume that the cluster under consideration, denoted $C_u$, has root at the origin $o$. Recall that $R_u(o)$ denotes the event that a cluster above $u$ has root at $o$. Therefore,
		\begin{equation}\label{eq:CSD_root_interpretation}
			\begin{split}
				\P(S_u=k)
				&= \P\bigl(|C_u|=k \,\big|\, \text{there exists a cluster $C_u$ above $u$ with root at $o$}\bigr)\\
				&= \P\bigl(|C_u(o)|=k\mid R_u(o)\bigr).
			\end{split}
		\end{equation}
		It then follows from
		\eqref{eq:w_d_2} that
		\[
		\P(S_u=k)
		=\frac{\P(R_u(o),|C_u(o)|=k)}{\P(R_u(o))}
		=\frac{w_k}{\rho_u}.
		\]
		By Assumption~\ref{ass:finite-cluster}, on $R_u(o)$ the cluster $C_u(o)$ has some finite size.  Hence
		\[
		R_u(o)=\bigsqcup_{j\ge1}
		\{R_u(o),\ |C_u(o)|=j\},
		\]
		together with \eqref{eq:w_d_2}, we obtain
		\[
		\rho_u = \P(R_u(o)) = \sum_{j=1}^\infty \P\bigl(R_u(o),\ |C_u(o)|=j\bigr) = \sum_{j=1}^\infty w_j,
		\]
		completing the proof.
	\end{proof}

\end{document}
